\Article{7}
Anstelle der Formel (3), die das Gewicht der plausibelsten Bestimmung ausdrückt, lassen sich mehrere andere angeben, deren Entwicklung sich lohnen wird.

Zunächst bemerken wir: Werden die Gleichungen (2) jeweils mit \(\frac ap,\ \frac{a'}{p'},\ \frac{a''}{p''}\) etc. multipliziert und addiert, so ergibt sich
\[\relax
[aa]x^0+[ab]y^0+[ac]z^0+\text{etc.}=\frac{aL}{p}+\frac{a'L'}{p'}+\frac{a''L''}{p''}+\text{etc.}
\]
Die linke Seite wird \(=0\); die rechte Seite bezeichnen wir entsprechend mit \([aL]\); wir haben daher
\[\relax
[aL]=0,\quad\text{und auf ganz gleiche Weise}\quad[bL]=0,\ [cL]=0\ \text{etc.}
\]
Multiplizieren wir ferner die Gleichungen (2) der Reihe nach mit \(\frac Lp,\ \frac{L'}{p'},\ \frac{L''}{p''}\) etc. und addieren sie, so finden wir
\[
\frac{lL}{p}+\frac{l'L'}{p'}+\frac{l''L''}{p''}+\text{etc.}=\frac{LL}{p}+\frac{L'L'}{p'}+\frac{L''L''}{p''}+\text{etc.}
\]
woraus wir einen \textit{zweiten} Ausdruck für das Gewicht erhalten,
\[
P=\frac{1}{\dfrac{lL}{p}+\dfrac{l'L'}{p'}+\dfrac{l''L''}{p''}+\text{etc.}}
\]
Multiplizieren wir schließlich die Gleichungen (2) der Reihe nach mit \(\frac lp,\ \frac{l'}{p'},\ \frac{l''}{p''}\) etc. und addieren sie, so gelangen wir zu einem \textit{dritten} Ausdruck für das Gewicht
\[
P=\frac{1}{[al]x^0+[bl]y^0+[cl]z^0+\text{etc.}+[ll]}
\]
wenn wir entsprechend den übrigen Bezeichnungen
\[
\frac{ll}{p}+\frac{l'l'}{p'}+\frac{l''l''}{p''}+\text{etc.}=[ll]
\]
setzen. Hieraus lässt sich mithilfe der Gleichungen (1) leicht zu einem \textit{vierten Ausdruck} übergehen, den wir folgendermaßen angeben:
\[
\begin{aligned}
\frac1P=[ll]&-[aa]x^0x^0-[bb]y^0y^0-[cc]z^0z^0-\text{etc.}\\
&-2[ab]x^0y^0-2[ac]x^0z^0-2[bc]y^0z^0-\text{etc.}
\end{aligned}
\]
\SourceLeaf{070}
\Article{8}
Die bisher entwickelte allgemeine Lösung ist vor allem dem Fall angepasst, dass \textit{eine} von den beobachteten Größen abhängige Unbekannte zu bestimmen ist. Sind jedoch die plausibelsten Werte mehrerer von denselben Beobachtungen abhängiger Unbekannter gesucht, oder ist noch ungewiss, welche Unbekannten vorzugsweise aus den Beobachtungen abzuleiten sind, so empfiehlt es sich, die Beobachtungen auf andere Weise vorzubereiten; diese Weise wollen wir nun entwickeln.

Wir betrachten die Größen \(x,y,z\) etc. als Unbestimmte, setzen
\[
\left.\begin{aligned}\relax
[aa]x+[ab]y+[ac]z+\text{etc.}&=\xi\\
[ab]x+[bb]y+[bc]z+\text{etc.}&=\eta\\
[ac]x+[bc]y+[cc]z+\text{etc.}&=\zeta
\end{aligned}\right\}\ (6)
\]
etc. und nehmen an, dass hieraus durch Elimination
\[
\left.\begin{aligned}\relax
[\alpha\alpha]\xi+[\alpha\beta]\eta+[\alpha\gamma]\zeta+\text{etc.}&=x\\
[\beta\alpha]\xi+[\beta\beta]\eta+[\beta\gamma]\zeta+\text{etc.}&=y\\
[\gamma\alpha]\xi+[\gamma\beta]\eta+[\gamma\gamma]\zeta+\text{etc.}&=z
\end{aligned}\right\}\ (7)
\]
etc. folgt.

Vor allem ist hier zu bemerken, dass symmetrisch angeordnete Koeffizienten notwendig gleich sind, beispielsweise
\[
\begin{aligned}\relax
[\beta\alpha]&=[\alpha\beta]\\
[\gamma\alpha]&=[\alpha\gamma]\\
[\gamma\beta]&=[\beta\gamma]\\
\text{etc.}&
\end{aligned}
\]
Dies folgt zwar unmittelbar aus der allgemeinen Theorie der Elimination bei linearen Gleichungen, wird aber unten auch ohne diese Theorie direkt bewiesen werden.

Wir haben daher
\[
\left.\begin{aligned}
x^0&=-[\alpha\alpha].[al]-[\alpha\beta].[bl]-[\alpha\gamma].[cl]-\text{etc.}\\
y^0&=-[\alpha\beta].[al]-[\beta\beta].[bl]-[\beta\gamma].[cl]-\text{etc.}\\
z^0&=-[\alpha\gamma].[al]-[\beta\gamma].[bl]-[\gamma\gamma].[cl]-\text{etc.}
\end{aligned}\right\}\ (8)
\]
\begin{center}etc.\end{center}
woraus wir, wenn wir
\SourceLeaf{071}
\[
\left.\begin{aligned}\relax
[\alpha\alpha]\mathfrak A+[\alpha\beta]\mathfrak B+[\alpha\gamma]\mathfrak C+\text{etc.}&=A\\
[\alpha\beta]\mathfrak A+[\beta\beta]\mathfrak B+[\beta\gamma]\mathfrak C+\text{etc.}&=B\\
[\alpha\gamma]\mathfrak A+[\beta\gamma]\mathfrak B+[\gamma\gamma]\mathfrak C+\text{etc.}&=C
\end{aligned}\right\}\ (9)
\]
etc. setzen,
\[
K=k-A[al]-B[bl]-C[cl]-\text{etc.}
\]
erhalten; oder wenn wir außerdem
\[
\left.\begin{aligned}
aA+bB+cC+\text{etc.}&=p\varepsilon\\
a'A+b'B+c'C+\text{etc.}&=p'\varepsilon'\\
a''A+b''B+c''C+\text{etc.}&=p''\varepsilon''
\end{aligned}\right\}\ (10)
\]
etc. setzen, gilt
\[
K=k-l\varepsilon-l'\varepsilon'-l''\varepsilon''-\text{etc.}\qquad(11)
\]
\clearpage
